\documentclass[10pt]{article}
\usepackage{icmlko}

\icmlmeta{IP 파운데이션 월드모델}{77}{기획서 한 장씩 넣어도 같은 자리였다}{마지막 두 누출을 막았다 --- 전도적 가정과 설계 결정의 시간 누출}

\begin{document}

\twocolumn[
\icmlheader

\begin{abstract}
노트 76이 시간을 갈랐지만 누출 둘을 남겼다. \textbf{첫째, 대상의 인자 공간을
미래 레코드 전체로 만들었다} --- 쉰아홉 건을 한꺼번에 보고 주성분을 뽑는 것은
제품이 쓰이는 방식이 아니다. 제품은 기획서 \emph{한 장}을 들고 온다. 규약 넷을
비교했다. \textbf{전부 같은 자리다} --- 한꺼번에 $+$0.5324, 하나 빼고(LOO)
$+$0.5313, 과거 공간에 미래 묶음 $+$0.5279, \textbf{완전 한 건씩 $+$0.5216}
{[}$+$0.246, $+$0.711{]}. 전도적 가정의 값이 \textbf{0.011}이다. 넷째가 제품
그대로다 --- 2025년 이전에 연 팝업 \textbf{16건}으로 공간을 만들고 새 기획을
\textbf{하나씩} 넣어 ``지금까지 연 것 중 상위 몇 퍼센트''를 낸다. \textbf{둘째,
배선과 축을 전체 데이터로 골랐다.} 2025년 이전만 보고 후보 79개를 다시 돌리니
전체 데이터의 $\Delta$와 \textbf{$r=+$0.867}로 맞고, 특히 애니 타깃 폭에서
태그 수가 과거 데이터로도 $-$0.0557이라 \textbf{2025년을 못 본 분석가도 같은
배선을 골랐다.} 다만 \emph{상위 순위}는 안 옮겨간다 --- 과거만 데이터의 상위
다섯이 전체에서 $+$0.0027$\sim-$0.0038로 전부 잡음이고, 그중 하나는 노트 67이
누출로 판정했던 \textbf{웹툰 회차 수}다. \textbf{표본을 자르면 누출이 다시
좋아 보인다.}
\end{abstract}
]

\begin{figure*}[t]
\centering
\includegraphics[width=\textwidth]{figs/one.pdf}
\caption{왼쪽 --- 규약 넷의 미래 팝업 $\rho$. 오른쪽 --- 과거만 데이터와 전체
데이터에서 본 배선 후보의 $\Delta$.}
\end{figure*}

\section{한 장씩 넣는다}

\begin{quote}\fontsize{9.5}{11.5}\selectfont
\textbf{한꺼번에} 대상 인자 공간을 평가 대상 전체로 만든다(노트 76까지)\\
\textbf{하나 빼고} 레코드 $i$를 뺀 나머지로 만들고 $i$를 투영한다\\
\textbf{과거 공간} $T$ 이전 대상 레코드로 만들고 미래를 묶어서 투영한다\\
\textbf{완전 한 건씩} 같은 공간에 미래를 \emph{하나씩} 넣는다 --- 제품 그대로
\end{quote}

\begin{table}[t]
\centering\tsize
\caption{시점 2025 · 미래 팝업 59건. 공간 적합용 과거는 16건.}
\begin{tabular}{lrl}
\toprule
규약 & $\rho$ & 95\% 구간 \\
\midrule
① 한꺼번에 & $+$0.5324 & --- \\
② 하나 빼고(LOO) & $+$0.5313 & --- \\
③ 과거 공간 $+$ 미래 묶음 & $+$0.5279 & $[+0.270, +0.722]$ \\
\textbf{④ 완전 한 건씩} & \textbf{$+$0.5216} & $[+0.246, +0.711]$ \\
\bottomrule
\end{tabular}
\end{table}

\claimx{전도적 가정은 값을 안 치른다 --- 가장 느슨한 규약과 가장 엄격한 규약의
차이가 0.011이다. 기획서를 한 장씩 넣어도 묶어서 넣은 것과 같은 성적이 나온다.}
{네 규약 모두 팝업 하나 · 시점 하나에서 봤다. 구간 반폭이 0.23이라 0.011짜리
차이는 어차피 못 가른다 --- ``같다''는 것은 ``다르다는 증거가 없다''는 뜻이다.}

\paragraph{한 건은 줄을 못 세운다}
구현에서 함정을 하나 만났다. 앙상블이 순위 평균인데 레코드가 하나뿐이면
\texttt{rankdata} 가 항상 1을 낸다. 처음 짠 코드가 $\rho=-$0.006을 냈고 그것이
버그였다. \textbf{새 기획 하나를 적합 집합과 함께 투영해 그 안에서 자리를
잰다.} 제품이 내는 문장도 그것이다 --- ``지금까지 연 팝업 중 상위 몇 퍼센트인가.''

\section{2025년을 못 본 분석가}

배선과 축은 여전히 전체 데이터로 골랐다. 2025년 이전만 남기고 후보 79개를
다시 돌렸다(팝업은 과거가 16건뿐이라 이 분석에서 빠지고 일곱 도메인이다).

\begin{table}[t]
\centering\tsize
\caption{애니 타깃 폭 · 굿즈 규모 후보. 현행 대비 $\Delta$.}
\begin{tabular}{llrr}
\toprule
슬롯 & 후보 & 과거만 & 전체 \\
\midrule
타깃 폭 & 연령 등급(현행) & $\pm$0.0000 & $\pm$0.0000 \\
타깃 폭 & 태그 수 & $-$0.0557 & $-$0.0513 \\
타깃 폭 & 연령 등급(역) & $-$0.0982 & $-$0.0794 \\
타깃 폭 & 장르 수 & $-$0.1027 & $-$0.0781 \\
굿즈 규모 & 태그 수 & $-$0.0130 & $-$0.0175 \\
굿즈 규모 & 화수 & $-$0.0316 & $-$0.0322 \\
\bottomrule
\end{tabular}
\end{table}

노트 75의 결정 --- 태그 수를 버리고 연령 등급으로 --- 이 과거 데이터로도
그대로 나온다($-$0.0557). \textbf{2025년을 못 본 분석가도 같은 배선을 골랐다.}

\claimx{채택된 설계 결정은 시간 누출이 아니다. 후보 79개에서 과거만 $\Delta$와
전체 $\Delta$의 상관이 $+$0.867이고, 채택된 것들은 전부 $|\Delta|>0.03$인 큰
결정이라 부호가 바뀌지 않는다.}{$+$0.867은 큰 결정들이 끌어올린 값이다. 작은
결정에서는 상관이 훨씬 낮다 --- 바로 다음 절이 그것이다.}

\section{자르면 누출이 다시 좋아 보인다}

\begin{table}[t]
\centering\tsize
\caption{과거만 데이터에서 가장 좋아 보이는 다섯. 전체에서는 전부 잡음이다.}
\begin{tabular}{llrr}
\toprule
도메인 & 후보 & 과거만 & 전체 \\
\midrule
애니 & 더빙 & $+$0.0080 & $+$0.0027 \\
\textbf{웹툰} & \textbf{회차 수} & \textbf{$+$0.0068} & \textbf{$-$0.0038} \\
아이돌 & 인원 수 & $+$0.0034 & $-$0.0034 \\
웹툰 & 작가 수 & $+$0.0031 & $+$0.0001 \\
웹툰 & 작가 수 & $+$0.0027 & $+$0.0014 \\
\bottomrule
\end{tabular}
\end{table}

둘째 줄이 눈에 띈다. \textbf{웹툰 회차 수}는 노트 46이 ``연재 기간에 비례해
사후 누적된다''며 뺀 것이고 노트 67이 $-$0.1019로 명확히 나빠짐을 확인한
누출 후보다. 과거만 데이터에서는 $+$0.0068로 \emph{좋아 보인다}.

\claimx{표본을 시간으로 자르면 사후 누적 변수가 다시 좋아 보인다. 자른 표본에서는
누적 기간이 짧아 누출 신호가 약해지고, 남은 상관이 진짜처럼 보인다.}{과거만
$+$0.0068도 반폭 아래라 애초에 채택 대상이 아니다. ``속을 뻔했다''가 아니라
``속을 수 있는 구조다''까지만 말할 수 있다.}

\section{남은 것}

이것으로 알고 있던 누출을 전부 막았다.

\begin{table}[t]
\centering\tsize
\caption{누출 점검표.}
\begin{tabular}{ll}
\toprule
대상 라벨을 모형이 보는가 & 안 본다(설계상) \\
출처가 미래를 보는가 & 안 본다(노트 76) \\
대상 공간이 미래를 보는가 & 안 본다(규약 ④) \\
평가가 보정 표본을 쓰는가 & 안 쓴다(노트 70) \\
설계 결정이 미래를 보는가 & \textbf{본다 --- 그러나 결과가 같다} \\
\bottomrule
\end{tabular}
\end{table}

마지막 줄만 형식적으로 열려 있다. 배선을 전체 데이터로 골랐다는 사실 자체는
남아 있고, ``골랐어도 결과가 같다''는 것을 이 노트가 보였을 뿐이다. 완전히
막으려면 앞으로의 결정을 과거 데이터로만 내려야 한다.

\section{한계}

\begin{itemize}[leftmargin=1.1em,itemsep=0.15em]
\item 규약 넷의 차이가 0.011인데 구간 반폭이 0.23이다. 넷이 같다는 것은
      ``다르다는 증거가 없다''는 뜻이다.
\item 과거 공간을 팝업 \textbf{16건}으로 만든다. 다섯 축에 16건이라 주성분이
      몹시 불안정할 텐데 성적이 안 떨어진 것이 오히려 이상하다.
\item 시간 분할 배선 탐색에서 팝업이 빠졌다(과거 16건 $<$ 문턱 40).
      팝업 배선 결정은 이 검사를 못 받았다.
\item 누출 점검표의 마지막 줄은 여전히 열려 있다.
\item 애니 후보 하나가 과거만 0.000, 전체 $-$0.115로 크게 어긋난다. 안 봤다.
\end{itemize}

\section{다음}

\begin{enumerate}[leftmargin=1.3em,itemsep=0.15em]
\item 앞으로의 설계 결정은 \textbf{과거 데이터로만} 내리고 미래로 검증한다.
      규약을 코드에 넣는다.
\item 아홉째 도메인 --- 배선이 닫혔으므로(노트 76) 남은 지렛대는 이것뿐이다.
\item 팝업 과거 표본이 16건이라 규약 ④가 불안정하다. 시점을 옮겨 가며
      민감도를 본다.
\item 애니가 시간을 자를 때 좋아지는 이유(노트 76의 남은 숙제).
\end{enumerate}

\vspace{0.6em}
\hrule height 0.3pt
\vspace{0.5em}
{\footnotesize
\textbf{참고문헌}\\[0.3em]
Kaufman, S., Rosset, S., Perlich, C. (2012). Leakage in data mining.
\emph{ACM TKDD}, 6(4), 1--21.\\[0.15em]
Bergmeir, C., Benitez, J. M. (2012). On the use of cross-validation for time
series predictor evaluation. \emph{Information Sciences}, 191, 192--213.\\[0.15em]
Cawley, G. C., Talbot, N. L. C. (2010). On over-fitting in model selection and
subsequent selection bias in performance evaluation. \emph{JMLR}, 11, 2079--2107.\\[0.15em]
Vapnik, V. (1998). \emph{Statistical Learning Theory}. Wiley.\\[0.15em]
Efron, B., Tibshirani, R. J. (1993). \emph{An Introduction to the Bootstrap}.
Chapman \& Hall.
}

\end{document}
